\documentclass{beamer}
% COMSC-3013 Computer Architecture shared Beamer preamble.
% Week decks live one directory below weeks/ and input this file as:
%   % COMSC-3013 Computer Architecture shared Beamer preamble.
% Week decks live one directory below weeks/ and input this file as:
%   % COMSC-3013 Computer Architecture shared Beamer preamble.
% Week decks live one directory below weeks/ and input this file as:
%   \input{../_shared/beamer-preamble.tex}
% Keep the course self-contained. Change visual/build conventions here,
% not by forking one-off week themes.

\usetheme{Boadilla}
\usecolortheme{seahorse}
\setbeamertemplate{navigation symbols}{}
\setbeamertemplate{footline}[frame number]

\usepackage{amsmath}
\usepackage{booktabs}
\usepackage{graphicx}
\usepackage{listings}
\usepackage{pgfpages}

\lstset{
  basicstyle=\ttfamily\small,
  columns=fullflexible,
  keepspaces=true,
  showstringspaces=false,
  breaklines=true
}

% Speaker notes remain hidden in the ordinary student build. Define \shownotes
% before inputting the deck to create an instructor copy; see weeks/README.md.
\ifdefined\shownotes
  \setbeameroption{show notes on second screen=right}
\fi

\newcommand{\ArchTitleSlide}[4]{%
  % #1 week label, #2 title, #3 central question, #4 Wednesday prediction
  \begin{frame}
    \titlepage
  \end{frame}
  \begin{frame}{The machine question}
    \Large #3

    \vspace{1.2em}
    \normalsize\textbf{Prediction to carry into Wednesday:} #4
  \end{frame}
}

% Keep the course self-contained. Change visual/build conventions here,
% not by forking one-off week themes.

\usetheme{Boadilla}
\usecolortheme{seahorse}
\setbeamertemplate{navigation symbols}{}
\setbeamertemplate{footline}[frame number]

\usepackage{amsmath}
\usepackage{booktabs}
\usepackage{graphicx}
\usepackage{listings}
\usepackage{pgfpages}

\lstset{
  basicstyle=\ttfamily\small,
  columns=fullflexible,
  keepspaces=true,
  showstringspaces=false,
  breaklines=true
}

% Speaker notes remain hidden in the ordinary student build. Define \shownotes
% before inputting the deck to create an instructor copy; see weeks/README.md.
\ifdefined\shownotes
  \setbeameroption{show notes on second screen=right}
\fi

\newcommand{\ArchTitleSlide}[4]{%
  % #1 week label, #2 title, #3 central question, #4 Wednesday prediction
  \begin{frame}
    \titlepage
  \end{frame}
  \begin{frame}{The machine question}
    \Large #3

    \vspace{1.2em}
    \normalsize\textbf{Prediction to carry into Wednesday:} #4
  \end{frame}
}

% Keep the course self-contained. Change visual/build conventions here,
% not by forking one-off week themes.

\usetheme{Boadilla}
\usecolortheme{seahorse}
\setbeamertemplate{navigation symbols}{}
\setbeamertemplate{footline}[frame number]

\usepackage{amsmath}
\usepackage{booktabs}
\usepackage{graphicx}
\usepackage{listings}
\usepackage{pgfpages}

\lstset{
  basicstyle=\ttfamily\small,
  columns=fullflexible,
  keepspaces=true,
  showstringspaces=false,
  breaklines=true
}

% Speaker notes remain hidden in the ordinary student build. Define \shownotes
% before inputting the deck to create an instructor copy; see weeks/README.md.
\ifdefined\shownotes
  \setbeameroption{show notes on second screen=right}
\fi

\newcommand{\ArchTitleSlide}[4]{%
  % #1 week label, #2 title, #3 central question, #4 Wednesday prediction
  \begin{frame}
    \titlepage
  \end{frame}
  \begin{frame}{The machine question}
    \Large #3

    \vspace{1.2em}
    \normalsize\textbf{Prediction to carry into Wednesday:} #4
  \end{frame}
}

\title{Race Conditions in the Wild}
\subtitle{Agents, Git, networks, and shared mutable state}
\author{COMSC-3013 Computer Architecture}
\date{}

\begin{document}

\begin{frame}\titlepage\end{frame}

\begin{frame}{The incident}
\Large We added a second Foreman because two workers should get more done.

\vspace{1em}
\normalsize
\begin{itemize}
  \item Lead Foreman pulled from the top of the task queue.
  \item Assistant Foreman pulled from the bottom.
  \item Their tasks were logically different.
  \item They still touched the same mutable Git checkout.
\end{itemize}

\vspace{0.7em}
\centering
\textbf{Nothing collided in the plan. The state collided anyway.}
\note{Tell the real story: the queue design was sensible. The failure was shared mutable state, not bad task selection.}
\end{frame}

\begin{frame}{What is a race condition?}
\begin{block}{Working definition}
A race condition exists when the correctness or visible result depends on the timing or interleaving of operations that access shared state.
\end{block}

\begin{itemize}
  \item The workers can each be individually correct.
  \item The operations can each be individually valid.
  \item The combined result can still be wrong or unstable.
  \item The bug may disappear when you try to reproduce it.
\end{itemize}

\vspace{0.5em}
\textbf{Race conditions are about ordering plus shared state, not merely about threads.}
\end{frame}

\begin{frame}{The architecture pattern}
\begin{columns}[T]
\column{0.47\textwidth}
\begin{block}{Worker A}
read state\\
compute\\
write state
\end{block}

\column{0.47\textwidth}
\begin{block}{Worker B}
read state\\
compute\\
write state
\end{block}
\end{columns}

\vspace{1em}
\centering
\Large Shared mutable state

\vspace{0.6em}
\normalsize
If A and B can overlap, the system needs an ownership, synchronization, isolation, or message-passing rule.
\end{frame}

\begin{frame}{Agents do not make concurrency disappear}
\begin{itemize}
  \item Claude and Codex were different agents with different instructions.
  \item Separate prompts are not separate memory.
  \item Separate terminal sessions are not separate filesystems.
  \item Even separate machines can still race through a shared filesystem, database, API, or repository workflow.
\end{itemize}

\vspace{0.8em}
\begin{block}{Important distinction}
\textbf{Logical independence} of workers does not imply \textbf{resource independence} of the state they mutate.
\end{block}
\note{Connect this to cores: two cores can execute different instructions and still contend on the same cache line or memory location.}
\end{frame}

\begin{frame}{The network was not the isolation boundary}
\begin{columns}[T]
\column{0.45\textwidth}
\textbf{What looked separate}
\begin{itemize}
  \item different agent processes
  \item different task directions
  \item different terminal sessions
  \item potentially different network paths
\end{itemize}

\column{0.50\textwidth}
\textbf{What was actually shared}
\begin{itemize}
  \item one physical checkout
  \item one working tree
  \item one Git index
  \item one set of uncommitted files
  \item one stash/rebase state
\end{itemize}
\end{columns}

\vspace{0.8em}
\centering
\textbf{Network separation can move workers apart without isolating the state that matters.}
\end{frame}

\begin{frame}{Git has more shared state than the files you edit}
\small
\begin{tabular}{p{0.30\textwidth}p{0.61\textwidth}}
\toprule
\textbf{Git surface} & \textbf{Why concurrent Foremen care} \\
\midrule
working tree & file contents can change underneath another process \\
index & staging is repository state, not a property of one shell \\
HEAD / branch & checkout and commit operations change context \\
stash / rebase / merge state & one process can alter recovery/integration state another process assumes \\
refs / origin & pushes can create ordinary divergence that must be reconciled \\
\bottomrule
\end{tabular}

\vspace{0.8em}
\textbf{``We edit different files'' is not a complete concurrency-control strategy.}
\end{frame}

\begin{frame}{A simplified failure timeline}
\small
\begin{enumerate}
  \item Lead Foreman reads the control repo and begins reconciliation.
  \item Assistant Foreman is already changing task/control-plane files in the same checkout.
  \item Lead observes concurrent edits and tries to preserve them while committing separate work.
  \item Files and Git state change between observations; a Foreman log edit appears and disappears during the race.
  \item Git eventually reports divergence and refuses unsafe integration steps.
  \item Lead Foreman stops mutation rather than guessing which state should win.
\end{enumerate}

\vspace{0.6em}
\begin{block}{The safety success}
The system did not avoid the race. It \textbf{detected enough inconsistency to stop before turning a transient interleaving into silent corruption}.
\end{block}
\end{frame}

\begin{frame}{Why top-of-queue vs bottom-of-queue was not enough}
\begin{columns}[T]
\column{0.47\textwidth}
\begin{block}{Task-level partitioning}
Lead: top\\
Assistant: bottom\\
Different intended jobs
\end{block}

\column{0.47\textwidth}
\begin{block}{Resource-level reality}
Same checkout\\
Same index\\
Same control files
\end{block}
\end{columns}

\vspace{1em}
\centering
\Large Partitioning \textit{work} does not partition \textit{state}.

\vspace{0.6em}
\normalsize
This is the same reason two cores can execute different functions yet still race on one shared variable.
\end{frame}

\begin{frame}{The repair: isolation plus ownership}
\begin{columns}[T]
\column{0.47\textwidth}
\textbf{Filesystem isolation}
\begin{itemize}
  \item one Git worktree per Foreman seat
  \item separate working tree and per-worktree Git state
  \item shared repository history without shared mutable checkout state
\end{itemize}

\column{0.47\textwidth}
\textbf{Single-writer ownership}
\begin{itemize}
  \item Lead owns canonical queue/control files
  \item Assistant writes completion receipts
  \item Chat sessions submit proposals
  \item integration happens through Git, not simultaneous editing
\end{itemize}
\end{columns}
\note{Emphasize both parts. Isolation alone reduces accidental collision; ownership makes authority explicit.}
\end{frame}

\begin{frame}{Why worktrees are an Architecture idea}
\begin{block}{Analogy, not identity}
A Git worktree gives each Foreman a separate mutable execution context while still sharing durable repository objects/history.
\end{block}

\begin{itemize}
  \item Similar goal to separate process address spaces: isolate mutable state.
  \item Similar goal to per-core/private cache state: reduce unnecessary sharing.
  \item Git integration becomes an explicit synchronization point.
  \item Conflicts become visible events instead of invisible timing accidents.
\end{itemize}

\vspace{0.5em}
\textbf{Good systems make dangerous sharing explicit.}
\end{frame}

\begin{frame}{Three ways to tame shared state}
\begin{enumerate}
  \item \textbf{Synchronize it} - locks, transactions, atomic operations, serialized critical sections.
  \item \textbf{Partition it} - give each worker its own state and merge at explicit boundaries.
  \item \textbf{Own it} - designate one writer; everyone else sends messages/receipts/proposals.
\end{enumerate}

\vspace{1em}
Our Foreman redesign uses mostly \textbf{partition + ownership}, with Git as an explicit reconciliation mechanism.
\end{frame}

\begin{frame}{The same bug wears many costumes}
\begin{itemize}
  \item two CPU cores update one counter;
  \item two threads mutate one data structure;
  \item two services update one database row;
  \item two deployment jobs rewrite one environment;
  \item two agents edit one repository checkout;
  \item two humans overwrite the same shared document.
\end{itemize}

\vspace{0.8em}
\centering
\Large Different costume. Same question:

\vspace{0.4em}
\textbf{Who owns this state, and what orders the writes?}
\end{frame}

\begin{frame}{Closing rule}
\Large
More workers is a hypothesis.

\vspace{0.8em}
Parallelism only helps when the useful work outruns the costs of communication, synchronization, contention, and recovery.

\vspace{1em}
\normalsize
\begin{block}{Our incident in one sentence}
The agents were not fighting over the task queue; they were racing over the mutable machinery used to represent and coordinate that queue.
\end{block}
\end{frame}

\end{document}
